% GENERATED FILE. Edit canonical lesson frontmatter, then run npm run generate:books.
% canonical-entries: 39

\chapter*{Glossary}
\addcontentsline{toc}{chapter}{Glossary}
\markboth{Glossary}{Glossary}

This glossary collects every word and phrase taught in the book. Pronunciation appears when it differs from the written form; chapter numbers show where each entry is introduced.

\noindent\begin{minipage}[t]{\linewidth}
\raggedright
\textbf{\ru{брат}}\enspace\emph{brat}\par
\small brother — masculine, one soft sign away from a verb you already know, and this time it really is your English cousin\par
\footnotesize Introduced in Chapter 7.
\end{minipage}\par\medskip

\noindent\begin{minipage}[t]{\linewidth}
\raggedright
\textbf{\ru{брать}}\enspace\emph{brat'}\par
\small to take — English bear, and a partner from a completely different root\par
\footnotesize Introduced in Chapter 5.
\end{minipage}\par\medskip

\noindent\begin{minipage}[t]{\linewidth}
\raggedright
\textbf{\ru{быть}}\enspace\emph{byt'}\par
\small to be — and the one tense in which Russian refuses to say it\par
\footnotesize Introduced in Chapter 3.
\end{minipage}\par\medskip

\noindent\begin{minipage}[t]{\linewidth}
\raggedright
\textbf{\ru{чай}}\enspace\emph{chai}\par
\small tea — masculine, regular this time, and the overland twin of \ru{кофе}'s sea route\par
\footnotesize Introduced in Chapter 6.
\end{minipage}\par\medskip

\noindent\begin{minipage}[t]{\linewidth}
\raggedright
\textbf{\ru{читать}}\enspace\emph{chitát'}\par
\small to read — and why Russian has no separate way to say I am reading\par
\footnotesize Introduced in Chapter 4.
\end{minipage}\par\medskip

\noindent\begin{minipage}[t]{\linewidth}
\raggedright
\textbf{\ru{да}}\enspace\emph{da}\par
\small yes (da)\par
\footnotesize Introduced in Chapter 1.
\end{minipage}\par\medskip

\noindent\begin{minipage}[t]{\linewidth}
\raggedright
\textbf{\ru{друг}}\enspace\emph{drug}\par
\small friend (male, or unspecified) — masculine, and the reason you own \ru{ты} at all\par
\footnotesize Introduced in Chapter 7.
\end{minipage}\par\medskip

\noindent\begin{minipage}[t]{\linewidth}
\raggedright
\textbf{\ru{думать}}\enspace\emph{dúmat'}\par
\small to think — and the fact that shapes every Russian verb you will ever meet\par
\footnotesize Introduced in Chapter 4.
\end{minipage}\par\medskip

\noindent\begin{minipage}[t]{\linewidth}
\raggedright
\textbf{\ru{глаз}}\enspace\emph{glaz}\par
\small eye — masculine, and the word that quietly evicted Russian's inherited word for eye\par
\footnotesize Introduced in Chapter 9.
\end{minipage}\par\medskip

\noindent\begin{minipage}[t]{\linewidth}
\raggedright
\textbf{\ru{говорить}}\enspace\emph{govorít'}\par
\small to speak — one new letter, and a false friend worth killing early\par
\footnotesize Introduced in Chapter 3.
\end{minipage}\par\medskip

\noindent\begin{minipage}[t]{\linewidth}
\raggedright
\textbf{\ru{идти}}\enspace\emph{idtí}\par
\small to go — the one-way verb, and a past tense from a different root entirely\par
\footnotesize Introduced in Chapter 3.
\end{minipage}\par\medskip

\noindent\begin{minipage}[t]{\linewidth}
\raggedright
\textbf{\ru{как} \ru{вас} \ru{зовут}?}\enspace\emph{kak vas zovút}\par
\small what's your name? — literally 'HOW do they call you?', not 'what'\par
\footnotesize Introduced in Chapter 2.
\end{minipage}\par\medskip

\noindent\begin{minipage}[t]{\linewidth}
\raggedright
\textbf{\ru{хлеб}}\enspace\emph{khleb}\par
\small bread — masculine, and not a Slavic word at all\par
\footnotesize Introduced in Chapter 6.
\end{minipage}\par\medskip

\noindent\begin{minipage}[t]{\linewidth}
\raggedright
\textbf{\ru{кофе}}\enspace\emph{kófe}\par
\small coffee — masculine, the one exception worth meeting on purpose\par
\footnotesize Introduced in Chapter 6.
\end{minipage}\par\medskip

\noindent\begin{minipage}[t]{\linewidth}
\raggedright
\textbf{\ru{любить}}\enspace\emph{lyubít'}\par
\small to love, and to like — English love itself, with a letter that appears from nowhere\par
\footnotesize Introduced in Chapter 5.
\end{minipage}\par\medskip

\noindent\begin{minipage}[t]{\linewidth}
\raggedright
\textbf{\ru{меня} \ru{зовут}…}\enspace\emph{menyá zovút}\par
\small my name is… — literally 'they call me', with no word for 'my' and no word for 'name'\par
\footnotesize Introduced in Chapter 2.
\end{minipage}\par\medskip

\noindent\begin{minipage}[t]{\linewidth}
\raggedright
\textbf{\ru{нос}}\enspace\emph{nos}\par
\small nose — masculine, and the same word as English nose and Latin nasus\par
\footnotesize Introduced in Chapter 9.
\end{minipage}\par\medskip

\noindent\begin{minipage}[t]{\linewidth}
\raggedright
\textbf{\ru{нет}}\enspace\emph{nyet}\par
\small no (nyet)\par
\footnotesize Introduced in Chapter 1.
\end{minipage}\par\medskip

\noindent\begin{minipage}[t]{\linewidth}
\raggedright
\textbf{\ru{очень} \ru{приятно}}\enspace\emph{óchen priyátno}\par
\small pleased to meet you — literally 'very pleasant', and a cousin of English FRIEND\par
\footnotesize Introduced in Chapter 2.
\end{minipage}\par\medskip

\noindent\begin{minipage}[t]{\linewidth}
\raggedright
\textbf{\ru{писать}}\enspace\emph{pisát'}\par
\small to write — one new letter, and the one stress mistake you must not make\par
\footnotesize Introduced in Chapter 4.
\end{minipage}\par\medskip

\noindent\begin{minipage}[t]{\linewidth}
\raggedright
\textbf{\ru{подруга}}\enspace\emph{podrúga}\par
\small friend (female) — feminine, and \ru{друг} with a prefix and an ending stitched on\par
\footnotesize Introduced in Chapter 7.
\end{minipage}\par\medskip

\noindent\begin{minipage}[t]{\linewidth}
\raggedright
\textbf{\ru{помогать}}\enspace\emph{pomogát'}\par
\small to help — built on can, which is English may and might\par
\footnotesize Introduced in Chapter 5.
\end{minipage}\par\medskip

\noindent\begin{minipage}[t]{\linewidth}
\raggedright
\textbf{\ru{понимать}}\enspace\emph{ponimát'}\par
\small to understand — literally to take hold of, which is what comprehend means too\par
\footnotesize Introduced in Chapter 4.
\end{minipage}\par\medskip

\noindent\begin{minipage}[t]{\linewidth}
\raggedright
\textbf{\ru{пожалуйста}}\enspace\emph{pozhálusta}\par
\small please / you're welcome (pozhálusta)\par
\footnotesize Introduced in Chapter 1.
\end{minipage}\par\medskip

\noindent\begin{minipage}[t]{\linewidth}
\raggedright
\textbf{\ru{привет}}\enspace\emph{privét}\par
\small hi / hello (privét — informal)\par
\footnotesize Introduced in Chapter 1.
\end{minipage}\par\medskip

\noindent\begin{minipage}[t]{\linewidth}
\raggedright
\textbf{\ru{рот}}\enspace\emph{rot}\par
\small mouth — masculine, and this time the resemblance to English is not there to trust\par
\footnotesize Introduced in Chapter 9.
\end{minipage}\par\medskip

\noindent\begin{minipage}[t]{\linewidth}
\raggedright
\textbf{\ru{семья}}\enspace\emph{sem'yá}\par
\small family — feminine, and an unexpected cousin of the English word home\par
\footnotesize Introduced in Chapter 8.
\end{minipage}\par\medskip

\noindent\begin{minipage}[t]{\linewidth}
\raggedright
\textbf{\ru{сердце}}\enspace\emph{sérdtse}\par
\small heart — neuter, a silent letter, and one of the surest cognates in the entire Indo-European family\par
\footnotesize Introduced in Chapter 10.
\end{minipage}\par\medskip

\noindent\begin{minipage}[t]{\linewidth}
\raggedright
\textbf{\ru{сестра}}\enspace\emph{sestrá}\par
\small sister — feminine, brat's other half, and a cognate just as secure\par
\footnotesize Introduced in Chapter 7.
\end{minipage}\par\medskip

\noindent\begin{minipage}[t]{\linewidth}
\raggedright
\textbf{\ru{спасибо}}\enspace\emph{spasíbo}\par
\small thank you (spasíbo)\par
\footnotesize Introduced in Chapter 1.
\end{minipage}\par\medskip

\noindent\begin{minipage}[t]{\linewidth}
\raggedright
\textbf{\ru{спрашивать}}\enspace\emph{spráshivat'}\par
\small to ask a question — the verb behind the question you have asked since Chapter 2\par
\footnotesize Introduced in Chapter 5.
\end{minipage}\par\medskip

\noindent\begin{minipage}[t]{\linewidth}
\raggedright
\textbf{\ru{ты}, \ru{вы}}\enspace\emph{ty, vy}\par
\small you (informal) and you (formal) — the split English threw away\par
\footnotesize Introduced in Chapter 2.
\end{minipage}\par\medskip

\noindent\begin{minipage}[t]{\linewidth}
\raggedright
\textbf{\ru{ухо}}\enspace\emph{úkho}\par
\small ear — neuter, the first word Russian gives its third gender, and a word-for-word cousin of the English ear\par
\footnotesize Introduced in Chapter 9.
\end{minipage}\par\medskip

\noindent\begin{minipage}[t]{\linewidth}
\raggedright
\textbf{\ru{видеть}}\enspace\emph{vídet'}\par
\small to see — the root behind both vision and wisdom\par
\footnotesize Introduced in Chapter 3.
\end{minipage}\par\medskip

\noindent\begin{minipage}[t]{\linewidth}
\raggedright
\textbf{\ru{вода}}\enspace\emph{vodá}\par
\small water — feminine, and the first noun in this book, so it also teaches how every Russian noun announces itself\par
\footnotesize Introduced in Chapter 6.
\end{minipage}\par\medskip

\noindent\begin{minipage}[t]{\linewidth}
\raggedright
\textbf{\ru{я}}\enspace\emph{ya}\par
\small I — one letter, and one of the oldest words in Europe\par
\footnotesize Introduced in Chapter 2.
\end{minipage}\par\medskip

\noindent\begin{minipage}[t]{\linewidth}
\raggedright
\textbf{\ru{здравствуйте}}\enspace\emph{zdrávstvuyte}\par
\small hello (zdrávstvuyte — formal / polite)\par
\footnotesize Introduced in Chapter 1.
\end{minipage}\par\medskip

\noindent\begin{minipage}[t]{\linewidth}
\raggedright
\textbf{\ru{жить}}\enspace\emph{zhit'}\par
\small to live — and the English word that used to mean 'alive'\par
\footnotesize Introduced in Chapter 3.
\end{minipage}\par\medskip

\noindent\begin{minipage}[t]{\linewidth}
\raggedright
\textbf{\ru{знать}}\enspace\emph{znat'}\par
\small to know — the same root as English know, and the sentence you will use most\par
\footnotesize Introduced in Chapter 3.
\end{minipage}\par\medskip
